\documentclass[tikz,border=6pt]{standalone}
\usepackage{ctex}
\usepackage{amsmath}
\usetikzlibrary{calc, arrows.meta}

\begin{document}
\begin{tikzpicture}[scale=1.1, thick]

% 抛物线 y^2 = 4x，焦点 F(1,0)
% 焦点弦：过 F(1,0) 的弦 AB
% 调和关系示意：1/|AF| + 1/|BF| = 1/p（p 为半通径）

% 坐标轴
\draw[->, gray, thin] (-0.5, 0) -- (5.0, 0) node[right, font=\footnotesize] {$x$};
\draw[->, gray, thin] (0, -3.2) -- (0, 3.2) node[above, font=\footnotesize] {$y$};
\node[font=\footnotesize, below left] at (0,0) {$O$};

% 准线 x = -1
\draw[dashed, gray] (-1, -3.0) -- (-1, 3.0);
\node[font=\footnotesize, above] at (-1, 3.0) {准线 $x{=}{-1}$};

% 抛物线 y^2 = 4x（绘制若干点）
\draw[black, thick, domain=-3.0:3.0, samples=80, variable=\t]
    plot ({(\t)^2/4}, \t);

% 焦点 F(1,0)
\coordinate (F) at (1, 0);
\fill[red] (F) circle [radius=2pt];
\node[font=\footnotesize, below right] at (F) {$F(1,0)$};

% 焦点弦：斜率 k=1.2，直线 y = 1.2(x-1)
% 代入 y^2 = 4x: 1.44(x-1)^2 = 4x
% 1.44x^2 - 2.88x + 1.44 = 4x
% 1.44x^2 - 6.88x + 1.44 = 0
% x = (6.88 ± sqrt(47.334 - 8.2944)) / 2.88
% x = (6.88 ± sqrt(39.04)) / 2.88
% sqrt(39.04) ≈ 6.248
% x1 = (6.88 + 6.248)/2.88 ≈ 4.558
% x2 = (6.88 - 6.248)/2.88 ≈ 0.219
% y1 = 1.2*(4.558-1) = 4.270  (太大，调小斜率)
% 改用 k=0.7: y=0.7(x-1)
% 0.49(x-1)^2 = 4x
% 0.49x^2 - 0.98x + 0.49 = 4x
% 0.49x^2 - 4.98x + 0.49 = 0
% x = (4.98 ± sqrt(24.8 - 0.9604))/0.98 = (4.98 ± sqrt(23.84))/0.98
% sqrt(23.84) ≈ 4.883
% x1 = (4.98 + 4.883)/0.98 ≈ 10.06 (出界)
% 改用 k=1.5
% 2.25(x-1)^2 = 4x
% 2.25x^2 - 4.5x + 2.25 = 4x
% 2.25x^2 - 8.5x + 2.25 = 0
% x = (8.5 ± sqrt(72.25 - 20.25))/4.5 = (8.5 ± sqrt(52))/4.5
% sqrt(52) ≈ 7.211
% x1 = (8.5 + 7.211)/4.5 ≈ 3.491, y1 = 1.5*(3.491-1) = 3.737 (太高)
% 试 k=2.5: 6.25(x-1)^2 = 4x
% 6.25x^2 - 12.5x + 6.25 = 4x
% 6.25x^2 - 16.5x + 6.25 = 0
% x = (16.5 ± sqrt(272.25 - 156.25))/12.5 = (16.5 ± sqrt(116))/12.5
% sqrt(116)≈10.77
% x1=(16.5+10.77)/12.5≈2.182, y1=2.5*(2.182-1)=2.955
% x2=(16.5-10.77)/12.5≈0.458, y2=2.5*(0.458-1)=-1.355
\coordinate (A) at (2.182, 2.955);
\coordinate (B) at (0.458, -1.355);

% 弦 AB
\draw[blue, thick] (A) -- (B);
\fill[black] (A) circle [radius=2pt];
\fill[black] (B) circle [radius=2pt];
\node[font=\footnotesize, above right] at (A) {$A(x_1,y_1)$};
\node[font=\footnotesize, below right] at (B) {$B(x_2,y_2)$};

% 焦半径 FA, FB
\draw[dashed, orange, thick] (F) -- (A);
\draw[dashed, orange, thick] (F) -- (B);

% 标注 |AF|, |BF|
\node[font=\footnotesize, orange, right] at ($(F)!0.5!(A)$) {$|AF|$};
\node[font=\footnotesize, orange, right] at ($(F)!0.5!(B)+(0.05,0)$) {$|BF|$};

% 准线到 A, B 的距离（虚线）
\draw[dashed, green!50!black] (-1, 2.955) -- (A);
\draw[dashed, green!50!black] (-1, -1.355) -- (B);
\fill[green!50!black] (-1, 2.955) circle [radius=1.5pt];
\fill[green!50!black] (-1, -1.355) circle [radius=1.5pt];
\node[font=\tiny, left, green!50!black] at (-1, 2.955) {$A'$};
\node[font=\tiny, left, green!50!black] at (-1, -1.355) {$B'$};
\node[font=\tiny, green!50!black] at (-0.4, 2.955) {\tiny $|AA'|=x_1{+}1$};

% 通径（竖直焦点弦 x=1）
\coordinate (L1) at (1, 2.0);
\coordinate (L2) at (1, -2.0);
\draw[red!40, dashed, thin] (L1) -- (L2);
\node[font=\tiny, red!60, right] at (1, 1.7) {通径};

% 关键公式框
\node[draw, rounded corners, font=\footnotesize, fill=yellow!10, align=left]
  at (3.5, -2.0)
  {调和关系：$\dfrac{1}{|AF|}+\dfrac{1}{|BF|}=\dfrac{1}{p}$\\[2pt]
   韦达：$x_1+x_2, \; x_1 x_2$\\[2pt]
   焦半径：$|AF|=x_1+1$（准线距）};

\end{tikzpicture}
\end{document}
